% !TEX root = ../main.tex 
% (this is to make TexShop know where your parent file is)

% This file isn't used in the template, it's just an example file. You can run it by compiling the testlayout.tex file. 

% make a table and give it a caption and label
% A table with a legend - it will align center and next to the text if possible

\begin{table}[htp]
\centering

\begin{talltblr}
    [
     caption = {Defined Serial Square showing interval distances},
     entry = {Serial Square},
     label = {tblr:serial-square} 
    ]
    {
     rows = {2em, rowsep = 2pt},
     columns = {2em, colsep = 2pt},
     hspan = minimal,
     row{1} = {bg = NiceGray, 0.2em, rowsep=0.5pt},
     width = 0.4\linewidth,
     colspec = {X[1,l] X[1,l]  X[1,l]  X[1,l] X[1,l] X[1,l] },
    }

 \SetCell[c=6]{c,f}{
    \begin{tblr}
        {
         hspan =  minimal,
         baseline = f,
         width = \linewidth,
         colspec = {X [1,r]  X[1,r] X[1,r] X[1,r] X[1,r]X[1,r]},
        }
        \SetCell{r}{\scriptsize ${-2}$} & 
        \SetCell{r}{\scriptsize ${+3}$} & 
        \SetCell{r}{\scriptsize ${-3}$} & 
        \SetCell{r}{\scriptsize ${+4}$} & 
        \SetCell{r}{\scriptsize ${+4}$} &
    \end{tblr}}
    & 2-6 & 3-6 & 4-6 & 5-6 & 6 
 \\ 

3 & 1 & 4 & 1 & 5 & 9 \\ 
1 & 2 & 3 & 3 & 4 & 4 \\ 
4 & 1 & 2 & 4 & 5 & 6 \\ 
1 & 2 & 3 & 4 & 6 & 6 \\ 
5 & 4 & 6 & 5 & 6 & 8 \\ 
9 & 2 & 1 & 5 & 1 & 3 \\ 



\end{talltblr}
% \caption[Serial Square]{Defined Serial Square showing interval distances}
% \label{table:serial-square}
\end{table}
%\caption{Serial Square}
%\label{serial-square}
